% EGA V, English translation by P. Blass, 1987 (cote 24): archive cote 24, site page 21 (PDF page 22).
% typescript. Draft compared: none.
% AI-transcribed by Claude Opus 5.5 (claude-opus-5-5, adaptive thinking, effort high), Message Batches API, batch ega24b (msgbatch_01RGJJeHdR96U16ZWFuUjtDx), 2026-11-13; not yet checked by eye.
\scanpage{24/21}
\mspage{21}

integral the same is true about $Y_\eta$.  Indeed geometrically integral
$=$ geometrically irreducible $+$ separable.

\uline{Corollary 3.3}

Let us assume that $k$ is algebraically closed and that for
every irreducible component $X_i$ of $X$ we have $\dim f(X_i) \geq 2$, and
suppose also that $X$ is $\sigma$-connected where $\sigma$ is the set of closed
subsets $Z$ of $X$ such that $\dim \overline{f(X)} = 0$ (i.e. for every such $Z$,
$X - Z$ is connected).  Under such conditions $Y_\eta$ is geometrically
connected over $K(\eta)$.

Indeed by a lemma that ought to appear in par. 6* with
Hartshorne's theorem, the hypothesis signifies that we can join
any two irreducible components $X'$ and $X''$ of $X$ by a chain of
irreducible components $X_0$ and $X',\ldots,X_n = X''$ such that two
consecutive ones have an intersection $\notin \sigma$ so that the inverse
images $X'_\eta$ and $X''_\eta$ are joined by a chain of $X_{i\eta}$ which are
geometrically connected over $k(\eta)$ by 3.1 and the intersection of
two consecutive ones being $\neq \emptyset$ by 2.3.

It follows (since $Y_\eta = X_\eta$ is the union of the $X_{i\eta}$, $X_i$ running
through the set of irreducible components of $X$) that $Y_\eta$ is
geometrically connected over $k(\eta)$ q.e.d.

\uline{Section 4}\quad \uline{Study of a Variable Hyperplane Section:  ``Sufficiently}
\uline{General'' Sections}.

We return to the general situation of Section 1, S being an
arbitrary prescheme.  Also, we suppose that $X$ is of finite
presentation over $S$.

\margin{* ASK A.G.}
